\section{Method Overview}
\label{sec:method_overview}

The pipeline is implemented in Python 3.10 using LangGraph infrastructure to build our iterative workflow. The program starts with prefiltering, then loops between the Actor and Critic agents before outputting the completed 3DGS scene.
The prefiltering process first uses Python ffmpeg with GPU acceleration to downsample and select frames. The selected frames are then passed to COLMAP, a Structure from Motion (SfM) technique, to create the point cloud for initializing the Gaussian Splatting process. COLMAP outputs several files, including images.bin, cameras.bin, and points3d.bin, which are essential files for training the Gaussian splatting.\\

\noindent The Actor agent starts with initial configuration values for the first agentic loop, which are passed into gsplat’s simple\_trainer.py to train the 3DGS reconstruction. The gsplat function also takes in the directory storing the COLMAP outputs and the output directory location. After the first agentic loop, the LLM is asked to update the hyperparameters given an initial prompt, the feedback from the Critic agent, and the current configuration values. The response is then parsed and written to config.yaml before gsplat is trained again using the new parameters.\\

\noindent The Critic agent evaluates the reconstruction based on several components, including the following metrics: PSNR, SSIM, and LPIPS. It retrieves a subset of the rendered frames, which are rasterized novel view images of the 3D reconstruction. The VLM is provided the metrics and image, and is prompted to provide a structured response for the following with an emphasis on floaters, blurring, and needle-shaped Gaussians: whether the reconstruction is accepted or rejected, critiques, and recommended hyperparameter changes. If the VLM accepts the reconstruction, the loop ends; otherwise, the response is passed to the Actor agent to update the training configuration in order to further refine the scene.\\

\noindent In order to test our pipeline, we utilized a sample video in 4K resolution from the DL3DV-10K dataset \cite{ling2024dl3dv}.
